\section{Background and Related Work}
\label{sec:background}

\subsection{Zoned Architecture and Atom Rearrangement}

A spatial light modulator (SLM) creates dense static traps in the storage zone for atoms not currently participating in two-qubit gates. The entanglement zone contains isolated trap pairs, termed gate-site pairs. Two atoms occupying a pair receive a global Rydberg pulse to execute a controlled-Z (CZ) gate~\cite{evered2023highfidelity,lin2025zac}. AOD mobile traps transport the participating atoms into the entanglement zone. In Fig.~\ref{fig:architecture-preliminaries}(a), $q_0$ and $q_1$ occupy one gate-site pair, while $q_2$ receives a local Raman rotation in a storage trap. Atom $q_3$ enters the entanglement zone from storage, illustrating inter-zone transport alongside gate execution.

AOD-assisted atom rearrangement comprises trap transfers between SLM and AOD traps and continuous transport of AOD rows and columns. Three structural constraints govern the parallel execution of these trajectories. Active rows and columns must preserve their relative order during motion. Same-row and same-column relationships must match between source and destination positions. During trap transfers, simultaneously activated rows and columns must not create unintended traps at stationary atoms, termed ghost spots or unintended intersections~\cite{bluvstein2022coherenttransport,stade2024abstractmodel,stade2025routingaware}. Two trajectories that exchange relative order violate the first constraint. Trajectories sharing a source row but ending on different rows violate the second. Two transfers forming an unintended intersection at stationary atom $q_2$ violate the third, as illustrated in Fig.~\ref{fig:architecture-preliminaries}(b)--(d). The compiler groups compatible trajectories from each layer into rearrangement batches. These batches execute sequentially, and the longest trajectory determines each batch's duration.

Addressable Raman beams implement single-qubit rotations, whereas global Rydberg pulses in the entanglement zone implement two-qubit CZ gates. Global pulses also excite atoms outside the target gates within the entanglement zone, causing idle-atom excitation errors~\cite{evered2023highfidelity,lin2025zac}. Inter-layer residency reduces trap transfers but exposes temporarily inactive atoms to more Rydberg pulses.

Inter-layer routing must account for the occupancy of moving and stationary atoms. Given source and destination positions, the compiler groups and orders trajectories, ensuring source positions are vacated before reoccupation~\cite{stade2024abstractmodel,stade2025routingaware}.

\begin{figure}[!t]
  \centering
  \resizebox{\columnwidth}{!}{% Shared vector-figure vocabulary for the Chinese IEEE manuscript.
% Every figure in this directory inputs this file, so the manuscript only
% needs the standard TikZ package and the libraries already loaded by
% paper_zh.tex.  Keep the guard: figures may be input more than once while
% editing or when a stand-alone smoke-test document is built.
\ifdefined\GALKFigureStyleLoaded\else
\def\GALKFigureStyleLoaded{1}

% Baseline-derived visual tokens: ICCAD uses a restrained blue/orange device
% palette, while ZAC uses conventional black axes and a small fixed plot set.
\definecolor{galkInk}{HTML}{231F20}
\definecolor{galkMuted}{HTML}{666666}
\definecolor{galkGrid}{HTML}{D6D6D6}
\definecolor{galkPaper}{HTML}{FFFFFF}
% Semantic colors are fixed across every figure: orange=entanglement zone,
% blue=storage zone, black=atom identity, red=invalidity, and blue dashes=future
% prediction.  Plot and search-group colours are deliberately independent of
% the two physical-zone colours.
\definecolor{galkStorage}{HTML}{0064BC}
\definecolor{galkStorageFill}{HTML}{E8F1F7}
\definecolor{galkEntangle}{HTML}{E37323}
\definecolor{galkEntangleFill}{HTML}{F9ECE3}
\definecolor{galkResident}{HTML}{485CB3}
\definecolor{galkResidentFill}{HTML}{EEF0F8}
\definecolor{galkGood}{HTML}{A4AF07}
\definecolor{galkBad}{HTML}{C84E45}
\definecolor{galkLook}{HTML}{0064BC}
\definecolor{galkData}{HTML}{315E8A}
\definecolor{galkTail}{HTML}{B44B27}
\definecolor{galkGroupOne}{HTML}{4C78A8}
\definecolor{galkGroupTwo}{HTML}{7A7A7A}
\definecolor{galkGroupThree}{HTML}{9C755F}

\tikzset{
  % Shared engineering-figure contract.  Physical snapshots must place an
  % orange entanglement band above a blue storage field; only their scale may
  % change.  Labels are aligned at the upper-left of both regions.
  galk/base/.style={
    font=\rmfamily\fontsize{9.0pt}{10.0pt}\selectfont,
    text=galkInk,
    line cap=butt,
    line join=miter,
    >=Stealth
  },
  galk/panel title/.style={
    font=\rmfamily\bfseries\fontsize{9.2pt}{10.2pt}\selectfont,
    text=galkInk,
    anchor=west
  },
  galk/label/.style={
    font=\rmfamily\fontsize{9.0pt}{10.0pt}\selectfont,
    text=galkInk
  },
  galk/small label/.style={
    font=\rmfamily\fontsize{8.7pt}{9.7pt}\selectfont,
    text=galkInk
  },
  galk/zone label/.style={
    font=\rmfamily\bfseries\fontsize{8.7pt}{9.7pt}\selectfont,
    text=galkInk,
    anchor=north west,
    inner sep=.5pt
  },
  galk/entanglement zone/.style={
    draw=galkEntangle,
    fill=galkEntangleFill,
    line width=.40pt
  },
  galk/storage zone/.style={
    draw=galkStorage,
    fill=galkStorageFill,
    line width=.40pt
  },
  galk/site/.style={
    circle,
    draw=galkInk!68,
    fill=white,
    line width=.36pt,
    minimum size=1.35mm,
    inner sep=0pt
  },
  galk/occupied atom/.style={
    circle,
    draw=galkInk,
    fill=galkInk,
    line width=.36pt,
    minimum size=1.55mm,
    inner sep=0pt
  },
  galk/target site/.style={
    circle,
    draw=galkInk!60,
    fill=galkInk!10,
    line width=.34pt,
    minimum size=1.55mm,
    inner sep=0pt
  },
  galk/q label/.style={
    font=\rmfamily\itshape\fontsize{8.5pt}{9.5pt}\selectfont,
    text=galkInk,
    inner sep=.25pt
  },
  galk/rydberg pair/.style={
    draw=galkInk!62,
    densely dashed,
    line width=.36pt
  },
  galk/accepted path/.style={
    -{Stealth[length=.95mm,width=.62mm]},
    draw=galkInk,
    line width=.48pt
  },
  galk/alternative path/.style={
    -{Stealth[length=.95mm,width=.62mm]},
    draw=galkMuted,
    densely dashed,
    line width=.48pt
  },
  galk/future path/.style={
    -{Stealth[length=.95mm,width=.62mm]},
    draw=galkLook,
    densely dashed,
    line width=.48pt
  },
  galk/invalid path/.style={
    -{Stealth[length=.95mm,width=.62mm]},
    draw=galkBad,
    densely dashed,
    line width=.50pt
  }
}
\fi

\begin{tikzpicture}[
  galk/base,x=1cm,y=1cm,
  trap/.style={galk/site},
  atom/.style={galk/occupied atom},
  fixed/.style={circle,draw=galkInk!65,line width=.35pt,fill=galkInk!25,
    minimum size=1.55mm,inner sep=0pt},
  motion/.style={galk/accepted path},
  guide/.style={draw=galkStorage!45,densely dashed,line width=.40pt},
  inactive guide/.style={draw=galkGrid,densely dashed,line width=.35pt},
  panel title/.style={galk/panel title},
  relation/.style={font=\rmfamily\fontsize{9.3pt}{10.3pt}\selectfont},
  case label/.style={font=\rmfamily\bfseries\fontsize{9.2pt}{10.2pt}\selectfont},
  explanatory note/.style={font=\rmfamily\fontsize{9.2pt}{10.2pt}\selectfont},
  legend label/.style={font=\rmfamily\fontsize{9.2pt}{10.2pt}\selectfont},
  galk/q label/.append style={font=\rmfamily\itshape\fontsize{9.0pt}{10.0pt}\selectfont},
  constraint divider/.style={draw=galkInk!45,dash pattern=on 2pt off 2pt,line width=.45pt}
]
  \path[use as bounding box] (0,-.65) rectangle (17.80,7.20);

  % Device abstraction: rectangular zones and regular trap arrays.
  \node[panel title] at (.08,6.97) {(a) Zoned neutral-atom array};
  \draw[galk/entanglement zone] (.48,4.48) rectangle (7.40,6.55);
  \draw[galk/storage zone] (.48,1.02) rectangle (7.40,4.10);
  \node[galk/zone label] at (.66,6.39) {entanglement zone};
  \node[galk/zone label] at (.66,3.94) {storage zone};

  % Five equally spaced Rydberg sites, with two traps per site.
  \foreach \cx in {1.18,2.50,3.82,5.14,6.46} {
    \draw[galk/rydberg pair] (\cx,5.12) ellipse (.35 and .25);
    \node[trap] at (\cx-.14,5.12) {};
    \node[trap] at (\cx+.14,5.12) {};
  }
  \node[atom] (qzero) at (3.68,5.12) {};
  \node[atom] (qone) at (3.96,5.12) {};
  \draw[galkInk,line width=.48pt] (qzero)--(qone);
  \node[galk/q label,anchor=north east] at (3.62,4.81) {$q_0$};
  \node[galk/q label,anchor=north west] at (4.02,4.81) {$q_1$};
  \node[anchor=south] (rydberg) at (3.91,5.66) {Rydberg pulse (2Q)};
  \draw[motion] (rydberg.south)--(qone.north);

  % Eleven columns and five rows in the storage zone.
  \foreach \x in {.92,1.51,2.10,2.69,3.28,3.87,4.46,5.05,5.64,6.23,6.82} {
    \foreach \y in {1.42,1.85,2.28,2.71,3.14} {
      \node[trap] at (\x,\y) {};
    }
  }
  \node[atom] (qtwo) at (2.69,3.14) {};
  \node[galk/q label,anchor=north east] at (2.62,3.03) {$q_2$};
  \node[anchor=south] (raman) at (3.87,3.46) {Raman pulse (1Q)};
  \draw[motion] (raman.south west)--(qtwo.north east);

  % A hollow destination and one arrow denote a single transported atom.
  \node[atom] (qthree) at (5.64,2.28) {};
  \node[galk/q label,anchor=north east] at (5.57,2.16) {$q_3$};
  \node[trap] (qthreeDest) at (6.32,5.12) {};
  \draw[motion] (qthree.north east)
    .. controls (6.02,3.12) and (6.23,4.23) .. (qthreeDest.south);
  \node[galk/q label,anchor=north east] at (6.26,4.81) {$q_3$};
  \node[anchor=east] at (5.86,4.28) {AOD transport};
  \draw[galkInk,line width=.38pt] (5.88,4.28)--(6.14,4.28);

  \draw[-{Stealth[length=.9mm,width=.6mm]},line width=.42pt]
    (.26,.84)--(7.53,.84) node[below left,inner sep=1pt] {$x$};
  \draw[-{Stealth[length=.9mm,width=.6mm]},line width=.42pt]
    (.26,.84)--(.26,6.63);
  \node[anchor=center,inner sep=0pt] at (.12,6.48) {$y$};
  \draw[{Stealth[length=.8mm,width=.5mm]}-{Stealth[length=.8mm,width=.5mm]},
    line width=.38pt] (1.18,4.70)--(2.50,4.70);
  \node[fill=galkEntangleFill,inner sep=.7pt] at (1.84,4.70) {$d_{\rm site}$};
  \node[anchor=west,text=galkMuted,explanatory note]
    at (.54,.62) {Schematic; not to scale};

  % Strict ordering and row/column membership are one relation constraint.
  \node[panel title] at (8.02,6.97) {(b) Row/column relation constraint};
  \node[anchor=west] at (8.02,6.57) {Preserve $<$, $=$ and $>$ on both axes};
  \node[case label] at (9.48,6.20) {Valid};
  \node[case label] at (12.65,6.20) {Order reversal};
  \node[case label] at (15.91,6.20) {Row merging};
  \draw[galkGrid,line width=.38pt] (11.01,3.73)--(11.01,6.30);
  \draw[galkGrid,line width=.38pt] (14.27,3.73)--(14.27,6.30);

  % Valid: the common row and the order of two columns are preserved.
  \draw[guide] (8.20,4.72)--(10.79,4.72);
  \draw[guide] (8.20,5.73)--(10.79,5.73);
  \node[atom] (v0s) at (8.47,4.72) {};
  \node[atom] (v1s) at (9.66,4.72) {};
  \node[trap] (v0t) at (8.97,5.73) {};
  \node[trap] (v1t) at (10.16,5.73) {};
  \draw[motion] (v0s)--(v0t);
  \draw[motion] (v1s)--(v1t);
  \node[galk/q label,anchor=north] at (8.47,4.60) {$q_0$};
  \node[galk/q label,anchor=north] at (9.66,4.60) {$q_1$};
  \node[galk/q label,anchor=south] at (8.97,5.83) {$q_0$};
  \node[galk/q label,anchor=south] at (10.16,5.83) {$q_1$};
  \node[relation] at (9.48,4.14) {$x_0<x_1\;\to\;x'_0<x'_1$};
  \node[relation] at (9.48,3.68) {$y_0=y_1\;\to\;y'_0=y'_1$};

  % Invalid: column order reverses, while the common row remains unchanged.
  \draw[guide] (11.26,4.72)--(14.00,4.72);
  \draw[guide] (11.26,5.73)--(14.00,5.73);
  \node[atom] (o0s) at (11.69,4.72) {};
  \node[atom] (o1s) at (13.54,4.72) {};
  \node[trap] (o0t) at (13.54,5.73) {};
  \node[trap] (o1t) at (11.69,5.73) {};
  \draw[galk/invalid path] (o0s)--(o0t);
  \draw[galk/invalid path] (o1s)--(o1t);
  \node[galk/q label,anchor=north] at (11.69,4.60) {$q_0$};
  \node[galk/q label,anchor=north] at (13.54,4.60) {$q_1$};
  \node[galk/q label,anchor=south] at (13.54,5.83) {$q_0$};
  \node[galk/q label,anchor=south] at (11.69,5.83) {$q_1$};
  \node[relation,text=galkBad] at (12.65,4.14)
    {$x_0<x_1\;\to\;x'_0>x'_1$};
  \node[relation] at (12.65,3.68) {$y_0=y_1\;\to\;y'_0=y'_1$};

  % Invalid without a geometric crossing: distinct source rows merge.
  \draw[guide] (14.51,4.72)--(17.40,4.72);
  \draw[guide] (14.51,5.13)--(17.40,5.13);
  \draw[guide] (14.51,5.73)--(17.40,5.73);
  \node[atom] (m0s) at (14.92,4.72) {};
  \node[atom] (m1s) at (16.11,5.13) {};
  \node[trap] (m0t) at (15.51,5.73) {};
  \node[trap] (m1t) at (16.70,5.73) {};
  \draw[galk/invalid path] (m0s)--(m0t);
  \draw[galk/invalid path] (m1s)--(m1t);
  \node[galk/q label,anchor=north] at (14.92,4.60) {$q_0$};
  \node[galk/q label,anchor=north west] at (16.18,5.05) {$q_1$};
  \node[galk/q label,anchor=south] at (15.51,5.83) {$q_0$};
  \node[galk/q label,anchor=south] at (16.70,5.83) {$q_1$};
  \node[relation] at (15.91,4.14) {$x_0<x_1\;\to\;x'_0<x'_1$};
  \node[relation,text=galkBad] at (15.91,3.68)
    {$y_0<y_1\;\to\;y'_0=y'_1$};

  % Separate the two constraint classes without reducing the diagrams or text.
  \draw[constraint divider] (8.02,3.35)--(17.62,3.35);

  % The second constraint is the Cartesian product of activated beams.
  \begin{scope}[yshift=-.45cm]
  \node[panel title] at (8.02,3.33) {(c) Unintended-intersection constraint};
  \node[case label] at (9.48,2.93) {Simultaneous pickup};
  \node[case label] at (12.65,2.93) {Batch $B_1$: pick up $q_0$};
  \node[case label] at (15.91,2.93) {Batch $B_2$: pick up $q_1$};
  \draw[galkGrid,line width=.38pt] (11.01,1.04)--(11.01,3.02);

  % Identical coordinates; the same stationary atom q2 is shown each time.
  \foreach \dx in {0,3.26,6.52} {
    \begin{scope}[xshift=\dx cm]
      \foreach \x in {8.65,10.00} {
        \draw[inactive guide] (\x,1.12)--(\x,2.61);
      }
      \foreach \y in {1.36,2.36} {
        \draw[inactive guide] (8.32,\y)--(10.33,\y);
      }
      \node[trap] at (10.00,1.36) {};
      \node[fixed] at (8.65,2.36) {};
      \node[galk/q label,anchor=east] at (8.49,2.36) {$q_2$};
      \node[galk/q label,anchor=west] at (10.16,2.36) {$q_1$};
    \end{scope}
  }

  % Two rows and two columns also activate the occupied cross-intersection.
  \foreach \x in {8.65,10.00} {
    \draw[guide] (\x,1.12)--(\x,2.61);
  }
  \foreach \y in {1.36,2.36} {
    \draw[guide] (8.32,\y)--(10.33,\y);
  }
  \node[atom] at (8.65,1.36) {};
  \node[atom] at (10.00,2.36) {};
  \node[fixed] at (8.65,2.36) {};
  \node[galk/q label,anchor=north] at (8.65,1.24) {$q_0$};
  \draw[galkBad,line width=.7pt] (8.51,2.22)--(8.79,2.50);
  \draw[galkBad,line width=.7pt] (8.51,2.50)--(8.79,2.22);
  \node[explanatory note,text=galkBad,anchor=center]
    at (9.48,.67) {Unintended pickup of $q_2$};

  % One row--column pair is active in each of the two serialized pickups.
  \draw[guide] (11.91,1.12)--(11.91,2.61);
  \draw[guide] (11.58,1.36)--(13.59,1.36);
  \node[atom] at (11.91,1.36) {};
  \node[atom] at (13.26,2.36) {};
  \node[galk/q label,anchor=north] at (11.91,1.24) {$q_0$};
  \draw[motion] (14.00,1.85)--(14.50,1.85);
  \draw[guide] (16.52,1.12)--(16.52,2.61);
  \draw[guide] (14.84,2.36)--(16.85,2.36);
  \node[trap] at (15.17,1.36) {};
  \node[atom] at (16.52,2.36) {};
  \node[galk/q label,anchor=north] at (15.17,1.24) {vacated};
  \node[explanatory note,anchor=center] at (14.26,.67)
    {Sequential pickup leaves $q_2$ in the SLM};
  \end{scope}

  % A dedicated legend row does not obscure any physical component.
  \begin{scope}[yshift=-.45cm]
  \draw[galkGrid,line width=.38pt] (.48,.32)--(17.62,.32);
  \node[atom] at (.68,.02) {};
  \node[anchor=west,legend label] at (.88,.02) {atom};
  \node[fixed] at (2.09,.02) {};
  \node[anchor=west,legend label] at (2.29,.02) {stationary SLM atom};
  \node[trap] at (5.55,.02) {};
  \node[anchor=west,legend label] at (5.75,.02) {empty / target trap};
  \draw[guide] (8.84,.02)--(9.37,.02);
  \node[anchor=west,legend label] at (9.47,.02) {AOD beam position};
  \draw[motion] (12.39,.02)--(12.92,.02);
  \node[anchor=west,legend label] at (13.02,.02) {valid};
  \draw[galk/invalid path] (14.59,.02)--(15.12,.02);
  \node[anchor=west,legend label] at (15.22,.02) {invalid};
  \end{scope}
\end{tikzpicture}
}
  \caption{Zoned architecture and AOD rearrangement constraints. (a) Atom $q_2$ receives a Raman single-qubit gate in storage, while $q_0$ and $q_1$ receive a Rydberg two-qubit gate in the entanglement zone. Atom $q_3$ illustrates transport between the storage and entanglement zones. (b)--(d) illustrate the non-crossing, row/column preservation, and unintended-intersection constraints, respectively; the last concerns stationary atom $q_2$. The two loading operations in (d) require separate batches $B_1$ and $B_2$. Solid and red dashed lines denote feasible and incompatible movements, respectively. Zone dimensions and trap counts are schematic in this illustration.}
  \label{fig:architecture-preliminaries}
\end{figure}

We evaluate compilation results with the fidelity model for zoned neutral-atom architectures~\cite{lin2025zac}:
\begin{equation}
\begin{aligned}
\log F={}&N_1\log f_1+N_2\log f_2
+N_{\mathrm{exc}}\log f_{\mathrm{exc}}\\
&+N_{\mathrm{tran}}\log f_{\mathrm{tran}}
+\sum_{q\in\mathcal Q}\log\left(1-\frac{t_q}{T_2}\right).
\end{aligned}
\label{eq:fidelity}
\end{equation}
Here, $N_1$ and $N_2$ count single- and two-qubit gates, respectively. $N_{\mathrm{exc}}$ counts excitations of nonparticipating atoms in the entanglement zone, and $N_{\mathrm{tran}}$ counts per-atom trap transfers between SLM and AOD traps. The corresponding per-operation fidelity factors are $f_1$, $f_2$, $f_{\mathrm{exc}}$, and $f_{\mathrm{tran}}$. $t_q$ denotes the accumulated idle time of atom $q$, and $T_2$ denotes the coherence time. With fixed gate counts, compilation primarily improves fidelity by reducing trap transfers, idle-atom excitation, and coherence losses. Following Stade et al.'s routing evaluation~\cite{stade2025routingaware}, we report rearrangement batches and rearrangement time. The batch count measures the number of compatible trajectory groups requiring sequential execution; fewer batches indicate greater routing parallelism. Rearrangement time sums critical-trajectory durations across batches, measuring the total time required for atom transport.

\subsection{Neutral-Atom Compilation Methods}

Fixed-array compilation partitions circuits and places atoms to match hardware gate sets and connectivity. Geyser partitions circuits around native multiqubit gates, while GRAPHINE constructs application-specific Rydberg atom arrangements from weighted interaction graphs~\cite{patel2022geyser,patel2023graphine}. Schmid et al.'s hybrid mapping combines long-range gates and atom transport for neutral-atom devices with multiple connectivity mechanisms~\cite{schmid2024hybrid}.

Reconfigurable arrays let atoms move during circuit execution, requiring compilers to coordinate gates and transport. OLSQ-DPQA jointly solves gate scheduling, placement, and movement, while Enola separates two-qubit gate scheduling from parallel transport~\cite{tan2022mapping,tan2024olsqdpqa,tan2025enola}. Atomique and PARALLAX increase parallelism through multiple AOD arrays and SWAP-free movement, respectively~\cite{wang2024atomique,ludmir2024parallax}. Q-Pilot fixes data atoms in static traps and uses reusable mobile ancillary atoms for two-qubit interactions. Weaver supports retargeting across devices through an FPQA-specific intermediate representation and scalable optimization passes~\cite{wang2024qpilot,kirmemis2025weaver}. These methods focus on connectivity and parallelism in reconfigurable arrays. Zoned architectures also require choices among storage isolation, Rydberg exposure, and inter-zone transfers.

For zoned architectures, Stade et al. developed an abstract model and rearrangement algorithms~\cite{stade2024abstractmodel}. ZAC reuses entanglement-zone atoms through adjacent-layer matching and selects gate and storage sites through minimum-weight matching with a next-partner term~\cite{lin2025zac}. Stade et al.'s routing-aware placement, hereafter ICCAD/QMAP, uses A* to search gate and intermediate placements~\cite{stade2025routingaware}. Its cost accounts for compatible movement groups, each group's longest trajectory, and the next interaction partner. PowerMove jointly optimizes gate scheduling, placement, movement, and inter-zone communication. Mantra reduces atom round trips through program rewriting and preemptive scheduling~\cite{ruan2025powermove,jang2025mantra}. ZAC and ICCAD/QMAP represent reuse-aware matching and routing-aware A* placement, respectively. We select both as external baselines and compare end-to-end compilation under the same zoned architecture, inputs, and fidelity model.

\subsection{Inter-Layer Residency and Storage-Site Selection}

Inter-layer residency determines the zone an atom occupies between two interactions. Remaining in the entanglement zone reduces round-trip transfers, whereas returning to storage avoids Rydberg exposure in intermediate layers. The assigned storage site determines where later transport starts and whether the atom obstructs other paths. Residency, assigned storage sites, and subsequent gate sites jointly determine the physical cost of inter-layer rearrangement.

Consider the sequence of two nonadjacent interactions involving atom $q$. In $L_\ell$, it executes a CZ gate with $p$, then interacts with $r$ after two intermediate two-qubit layers, as shown in Fig.~\ref{fig:baseline-motivation}(a). Residency incurs two idle excitations, whereas returning requires a round-trip transfer sequence through storage; Fig.~\ref{fig:baseline-motivation}(b) compares these costs. The resident atom accumulates more idle excitations over longer layer intervals, while the assigned storage site affects subsequent rearrangement. If $q$ occupies $s_1$ on the direct transport path, it blocks the $u\to u'$ trajectory. Site $s_2$ lies outside this path and permits direct transport, as shown in Fig.~\ref{fig:baseline-motivation}(c). The occupancy conflict between stationary atom $q$ and the single $u\to u'$ trajectory requires another storage site or an intermediate stop to change the path.

\begin{figure*}[t]
  \centering
  \resizebox{\textwidth}{!}{% Shared vector-figure vocabulary for the Chinese IEEE manuscript.
% Every figure in this directory inputs this file, so the manuscript only
% needs the standard TikZ package and the libraries already loaded by
% paper_zh.tex.  Keep the guard: figures may be input more than once while
% editing or when a stand-alone smoke-test document is built.
\ifdefined\GALKFigureStyleLoaded\else
\def\GALKFigureStyleLoaded{1}

% Baseline-derived visual tokens: ICCAD uses a restrained blue/orange device
% palette, while ZAC uses conventional black axes and a small fixed plot set.
\definecolor{galkInk}{HTML}{231F20}
\definecolor{galkMuted}{HTML}{666666}
\definecolor{galkGrid}{HTML}{D6D6D6}
\definecolor{galkPaper}{HTML}{FFFFFF}
% Semantic colors are fixed across every figure: orange=entanglement zone,
% blue=storage zone, black=atom identity, red=invalidity, and blue dashes=future
% prediction.  Plot and search-group colours are deliberately independent of
% the two physical-zone colours.
\definecolor{galkStorage}{HTML}{0064BC}
\definecolor{galkStorageFill}{HTML}{E8F1F7}
\definecolor{galkEntangle}{HTML}{E37323}
\definecolor{galkEntangleFill}{HTML}{F9ECE3}
\definecolor{galkResident}{HTML}{485CB3}
\definecolor{galkResidentFill}{HTML}{EEF0F8}
\definecolor{galkGood}{HTML}{A4AF07}
\definecolor{galkBad}{HTML}{C84E45}
\definecolor{galkLook}{HTML}{0064BC}
\definecolor{galkData}{HTML}{315E8A}
\definecolor{galkTail}{HTML}{B44B27}
\definecolor{galkGroupOne}{HTML}{4C78A8}
\definecolor{galkGroupTwo}{HTML}{7A7A7A}
\definecolor{galkGroupThree}{HTML}{9C755F}

\tikzset{
  % Shared engineering-figure contract.  Physical snapshots must place an
  % orange entanglement band above a blue storage field; only their scale may
  % change.  Labels are aligned at the upper-left of both regions.
  galk/base/.style={
    font=\rmfamily\fontsize{9.0pt}{10.0pt}\selectfont,
    text=galkInk,
    line cap=butt,
    line join=miter,
    >=Stealth
  },
  galk/panel title/.style={
    font=\rmfamily\bfseries\fontsize{9.2pt}{10.2pt}\selectfont,
    text=galkInk,
    anchor=west
  },
  galk/label/.style={
    font=\rmfamily\fontsize{9.0pt}{10.0pt}\selectfont,
    text=galkInk
  },
  galk/small label/.style={
    font=\rmfamily\fontsize{8.7pt}{9.7pt}\selectfont,
    text=galkInk
  },
  galk/zone label/.style={
    font=\rmfamily\bfseries\fontsize{8.7pt}{9.7pt}\selectfont,
    text=galkInk,
    anchor=north west,
    inner sep=.5pt
  },
  galk/entanglement zone/.style={
    draw=galkEntangle,
    fill=galkEntangleFill,
    line width=.40pt
  },
  galk/storage zone/.style={
    draw=galkStorage,
    fill=galkStorageFill,
    line width=.40pt
  },
  galk/site/.style={
    circle,
    draw=galkInk!68,
    fill=white,
    line width=.36pt,
    minimum size=1.35mm,
    inner sep=0pt
  },
  galk/occupied atom/.style={
    circle,
    draw=galkInk,
    fill=galkInk,
    line width=.36pt,
    minimum size=1.55mm,
    inner sep=0pt
  },
  galk/target site/.style={
    circle,
    draw=galkInk!60,
    fill=galkInk!10,
    line width=.34pt,
    minimum size=1.55mm,
    inner sep=0pt
  },
  galk/q label/.style={
    font=\rmfamily\itshape\fontsize{8.5pt}{9.5pt}\selectfont,
    text=galkInk,
    inner sep=.25pt
  },
  galk/rydberg pair/.style={
    draw=galkInk!62,
    densely dashed,
    line width=.36pt
  },
  galk/accepted path/.style={
    -{Stealth[length=.95mm,width=.62mm]},
    draw=galkInk,
    line width=.48pt
  },
  galk/alternative path/.style={
    -{Stealth[length=.95mm,width=.62mm]},
    draw=galkMuted,
    densely dashed,
    line width=.48pt
  },
  galk/future path/.style={
    -{Stealth[length=.95mm,width=.62mm]},
    draw=galkLook,
    densely dashed,
    line width=.48pt
  },
  galk/invalid path/.style={
    -{Stealth[length=.95mm,width=.62mm]},
    draw=galkBad,
    densely dashed,
    line width=.50pt
  }
}
\fi

\begin{tikzpicture}[
  galk/base,x=1cm,y=1cm,
  trap/.style={galk/site},
  atom/.style={galk/occupied atom},
  transition/.style={galk/accepted path},
  panel title/.style={galk/panel title},
  state note/.style={font=\rmfamily\fontsize{8.7pt}{9.7pt}\selectfont},
  atom label/.style={galk/q label,font=\rmfamily\itshape\fontsize{8.8pt}{9.8pt}\selectfont},
  pulse/.style={draw=galkEntangle,densely dashed,line width=.42pt}
]
  \path[use as bounding box] (0,0) rectangle (17.80,5.85);

  % One regular optical-trap geometry is reused without deformation.
  \def\miniDevice#1#2{%
    \begin{scope}[shift={(#1,#2)}]
      \path[galk/entanglement zone] (0,.80) rectangle (1.12,1.30);
      \path[galk/storage zone] (0,0) rectangle (1.12,.70);
      \draw[galk/rydberg pair] (.27,1.05) ellipse (.15 and .13);
      \draw[galk/rydberg pair] (.85,1.05) ellipse (.15 and .13);
      \foreach \xx in {.18,.36,.76,.94}
        \node[trap] at (\xx,1.05) {};
      \foreach \xx in {.18,.56,.94} {
        \node[trap] at (\xx,.20) {};
        \node[trap] at (\xx,.50) {};
      }
    \end{scope}%
  }
  \def\gatePair#1#2#3#4{%
    \node[atom] at (#1,#2) {};
    \node[atom] at (#1+.18,#2) {};
    \draw[galkInk,line width=.42pt] (#1,#2)--(#1+.18,#2);
    \node[atom label,anchor=south] at (#1+.09,#2+.30) {$#3,#4$};
  }
  \def\siteAnnotation#1#2#3#4{%
    \node[state note] at (#1+.56,#2-.27) {$q$ at $#3$};
    \draw[galkInk!65,line width=.35pt]
      (#1+.38,#2-.09)--(#1+.38,#2+#4-.16)--(#1+.49,#2+#4-.08);
  }

  \draw[galkInk!24,line width=.38pt] (5.48,.35)--(5.48,5.58);
  \draw[galkInk!24,line width=.38pt] (11.15,.35)--(11.15,5.58);

  % (a) The interaction order determines when q is needed again.
  \node[panel title] at (.15,5.57) {(a) Non-adjacent interaction};
  \foreach \xx/\lab in {.20/$L_{\ell}$,1.49/$L_{\ell+1}$,
                       2.78/$L_{\ell+2}$,4.07/$L_{\ell+3}$} {
    \miniDevice{\xx}{3.08}
    \node[anchor=south] at (\xx+.56,4.92) {\lab};
  }
  \gatePair{.38}{4.13}{q}{p}
  \gatePair{2.25}{4.13}{u}{v}
  \gatePair{3.54}{4.13}{w}{x}
  \gatePair{4.25}{4.13}{q}{r}
  \foreach \xx in {1.67,2.96} {
    \node[atom] at (\xx,3.58) {};
    \node[atom label,anchor=center] at (\xx+.17,3.40) {$q$};
  }
  \node[state note,anchor=north] at (.76,2.88) {$\mathrm{CZ}(q,p)$};
  \node[state note,anchor=north] at (2.05,2.88) {$\mathrm{CZ}(u,v)$};
  \node[state note,anchor=north] at (3.34,2.88) {$\mathrm{CZ}(w,x)$};
  \node[state note,anchor=north] at (4.63,2.88) {$\mathrm{CZ}(q,r)$};
  \draw[galkInk!60,{Stealth[length=.75mm]}-{Stealth[length=.75mm]},line width=.40pt]
    (1.72,2.18)--(4.09,2.18);
  \node[fill=white,inner sep=1pt,state note] at (2.91,2.18) {two intervening layers};

  % The two-line zone key uses the same rectangular bands as the snapshots.
  \path[galk/entanglement zone] (.45,1.13) rectangle (.73,1.34);
  \node[anchor=west,state note] at (.88,1.235) {entanglement zone};
  \path[galk/storage zone] (.45,.60) rectangle (.73,.81);
  \node[anchor=west,state note] at (.88,.705) {storage zone};

  % (b) The gate sequence is identical; only q's inter-layer location differs.
  \node[panel title] at (5.68,5.57) {(b) Inter-layer residency};
  % Layer headers share the same top band as panels (a) and (c).
  % Place each residency description below its own row of snapshots.
  \node[anchor=west] at (5.73,2.80) {Entanglement-resident};
  \node[anchor=west] at (5.73,.53) {Storage-mediated};
  \foreach \xx/\lab in {5.73/$L_{\ell}$,7.04/$L_{\ell+1}$,
                       8.35/$L_{\ell+2}$,9.66/$L_{\ell+3}$} {
    \miniDevice{\xx}{3.08}
    \miniDevice{\xx}{.87}
    \node[anchor=south] at (\xx+.56,4.92) {\lab};
  }
  \gatePair{5.91}{4.13}{q}{p}
  \gatePair{7.80}{4.13}{u}{v}
  \gatePair{9.11}{4.13}{w}{x}
  \gatePair{9.84}{4.13}{q}{r}
  \foreach \xx in {7.22,8.53} {
    \node[atom] at (\xx,4.13) {};
    \node[atom label,anchor=south] at (\xx,4.43) {$q$};
    \draw[pulse] (\xx-.13,3.94)--(\xx-.13,4.32);
    \draw[pulse] (\xx+.13,3.94)--(\xx+.13,4.32);
  }
  \gatePair{5.91}{1.92}{q}{p}
  \gatePair{7.80}{1.92}{u}{v}
  \gatePair{9.11}{1.92}{w}{x}
  \gatePair{9.84}{1.92}{q}{r}
  \foreach \xx in {7.22,8.53} {
    \node[atom] at (\xx,1.37) {};
    \node[atom label,anchor=center] at (\xx+.17,1.19) {$q$};
  }

  % (c) s1 and s2 are real traps on the same regular storage lattice.
  \node[panel title] at (11.35,5.57) {(c) Storage-site dependence};
  \foreach \xx/\lab in {11.48/$L_{\ell}$,13.84/$L_{\ell+1}$,16.20/$L_{\ell+2}$} {
    \miniDevice{\xx}{3.08}
    \miniDevice{\xx}{.94}
    \node[anchor=south] at (\xx+.56,4.92) {\lab};
  }

  % The source state contains q in the entanglement zone and u in storage.
  \foreach \yy in {3.08,.94} {
    \node[atom] at (11.66,\yy+1.05) {};
    \node[atom label,anchor=south] at (11.66,\yy+1.35) {$q$};
    \node[atom] at (11.66,\yy+.50) {};
    \node[atom label,anchor=center] at (11.83,\yy+.32) {$u$};
  }
  \draw[transition] (12.78,3.66)--(13.66,3.66);
  \node[state note,anchor=south] at (13.22,3.80) {$q\to s_1$};
  \draw[transition] (12.78,1.52)--(13.66,1.52);
  \node[state note,anchor=south] at (13.22,1.66) {$q\to s_2$};

  % At s1, q occupies the middle trap of u's direct route.
  \foreach \xx in {13.84,16.20} {
    \node[atom] at (\xx+.18,3.58) {};
    \node[atom] at (\xx+.56,3.58) {};
    \draw[galk/invalid path] (\xx+.26,3.58)--(\xx+.86,3.58);
    \draw[galkBad,line width=.60pt] (\xx+.48,3.50)--(\xx+.64,3.66);
    \draw[galkBad,line width=.60pt] (\xx+.48,3.66)--(\xx+.64,3.50);
  }
  \siteAnnotation{13.84}{3.08}{s_1}{.50}
  \siteAnnotation{16.20}{3.08}{s_1}{.50}
  \draw[galk/invalid path] (15.14,3.66)--(16.02,3.66);
  \node[state note,anchor=south,text=galkBad] at (15.58,3.80) {Blocked};

  % At s2, the same single-atom route remains clear and u reaches its target.
  \node[atom] at (14.02,1.44) {};
  \node[atom] at (14.40,1.14) {};
  \draw[transition] (14.10,1.44)--(14.70,1.44);
  \node[atom] at (17.14,1.44) {};
  \node[atom] at (16.76,1.14) {};
  \node[atom label,anchor=center] at (16.97,1.26) {$u$};
  \siteAnnotation{13.84}{.94}{s_2}{.20}
  \siteAnnotation{16.20}{.94}{s_2}{.20}
  \draw[transition] (15.14,1.52)--(16.02,1.52);
  \node[state note,anchor=south] at (15.58,1.66) {Clear};
\end{tikzpicture}
}
  \caption{Residency and storage-site selection between nonadjacent interactions. (a) Atom $q$ participates in another CZ gate after two intermediate two-qubit layers. (b) Remaining in the entanglement zone reduces trap transfers, whereas returning to storage avoids idle excitation in intermediate layers. (c) When stationary atom $q$ occupies $s_1$ on the direct path, the $u\to u'$ trajectory is infeasible. Placing $q$ at $s_2$ outside this path makes the trajectory feasible. The former state requires another storage site or an intermediate stop. All states use the same zone geometry; dimensions and trap counts are not to scale.}
  \label{fig:baseline-motivation}
\end{figure*}

Choosing $s_1$ or $s_2$ requires the same current trap-transfer count but permits different subsequent trajectories. Inter-layer states must therefore record each atom's specific physical position. Evaluating future gate layers from each candidate's resulting state incorporates the storage site's transport consequences into the current decision.

\subsection{Our Contributions}
\label{sec:contributions}

GA-LK builds on reuse-aware compilation and routing-aware placement~\cite{lin2025zac,stade2025routingaware} with the following contributions:
\begin{itemize}
  \item Joint genetic search coordinates gate sites, residency, and storage sites for returning atoms. Each chromosome represents two-qubit gate sites and residency choices between nonadjacent layers. The $K$-best injective matchings assign distinct storage sites to returning atoms, coordinating all three placement decisions under one physical objective. For small encoded search spaces, the solver performs exact enumeration.
  \item Decayed multilayer physical look-ahead starts from each candidate's resulting state. It updates atom positions and clocks in subsequent two-qubit layer order. The evaluation accumulates costs from trap transfers, idle-atom excitation, transport, and relocation of blocking atoms. These costs guide layer-boundary selection through the current decision's effects on nonadjacent interactions.
  \item Trajectory conflict-graph coloring and occupancy-order checks construct executable rearrangement batches. The conflict graph encodes AOD row/column ordering, row/column preservation, and unintended-intersection constraints. DSATUR coloring~\cite{brelaz1979coloring} provides initial groups, followed by checks of source release and destination occupancy order. Executable batches and their durations allow joint search and multilayer look-ahead to evaluate routing parallelism and transport costs for candidate placements.
\end{itemize}
